\section{UI Design -- Mockup (Group Work)}

\subsection{Mục tiêu thiết kế giao diện}

UI phải giúp actor nhận biết trạng thái hiện tại, hành động có thể thực hiện và kết quả sau mỗi thao tác. Trạng thái từ dữ liệu IoT phải hiển thị thời điểm cập nhật; dữ liệu có thể đã cũ, tài nguyên \texttt{OutOfService} và thao tác bị từ chối phải được phân biệt rõ ràng.

\subsection{Kiến trúc thông tin và luồng điều hướng}

\ArtifactSection{INFORMATION ARCHITECTURE}{UI-IA-01}{Sơ đồ nhóm màn hình theo ba vai trò \texttt{ACT-STU}, \texttt{ACT-OPR}, \texttt{ACT-MNT}; thể hiện điểm vào và các khu vực chính.}{Kiến trúc thông tin của Smart E-Mobility Hub}

\ArtifactSection{SCREEN FLOW}{UI-FLOW-01}{Luồng màn hình cấp cao cho Sinh viên, Nhân viên vận hành và Nhân viên kỹ thuật; không thay thế activity diagram nghiệp vụ.}{Luồng điều hướng tổng quát}

\subsection{Screen Catalog}

\begin{longtable}{|L{34mm}|L{45mm}|L{24mm}|L{31mm}|}
  \caption{Danh mục màn hình dự kiến}\label{tab:phase2-screen-catalog}\\
  \hline
  \rowcolor{gray!15}
  \textbf{Screen ID} & \textbf{Tên màn hình} & \textbf{Actor} & \textbf{Use case liên quan}\\\hline
  \endfirsthead
  \hline
  \rowcolor{gray!15}
  \textbf{Screen ID} & \textbf{Tên màn hình} & \textbf{Actor} & \textbf{Use case liên quan}\\\hline
  \endhead
  SCR-STU-HUB-SEARCH & Tra cứu Hub & ACT-STU & UC-HUB-01\\\hline
  SCR-STU-HUB-DETAIL & Chi tiết Hub và tài nguyên & ACT-STU & UC-HUB-02\\\hline
  SCR-STU-HUB-ALT & Danh sách Hub thay thế & ACT-STU & UC-HUB-03\\\hline
  SCR-STU-RES-DETAIL & Xác nhận và theo dõi lượt đặt xe & ACT-STU & UC-RES-01, UC-RES-02\\\hline
  SCR-STU-TRIP & Nhận xe và chuyến đi đang hoạt động & ACT-STU & UC-TRIP-01, UC-TRIP-02\\\hline
  SCR-STU-PARK & Đặt chỗ và quản lý lượt đỗ xe & ACT-STU & UC-PARK-01--UC-PARK-03\\\hline
  SCR-STU-CHG & Đăng ký và theo dõi sạc & ACT-STU & UC-CHG-01, UC-CHG-02\\\hline
  SCR-OPR-CHG & Quản lý lịch sạc & ACT-OPR & UC-CHG-01--UC-CHG-03\\\hline
  SCR-OPR-DASHBOARD & Giám sát mạng lưới & ACT-OPR & UC-OPS-01\\\hline
  SCR-OPR-INCIDENT & Chi tiết và xử lý sự cố & ACT-OPR, ACT-MNT & UC-OPS-02\\\hline
  SCR-OPR-REDIST & Kế hoạch điều chuyển & ACT-OPR & UC-OPS-03\\\hline
  SCR-OPR-SIM-CONFIG & Cấu hình và chạy mô phỏng & ACT-OPR & UC-SIM-01\\\hline
  SCR-OPR-SIM-RESULT & Kết quả và khuyến nghị & ACT-OPR & UC-SIM-02\\\hline
\end{longtable}

\subsection{Trạng thái giao diện dùng chung}

Mỗi mockup cần thể hiện, khi có liên quan: trạng thái mặc định, loading, empty, dữ liệu có thể đã cũ, validation error, lỗi hệ thống, thao tác thành công và thao tác bị từ chối do business rule. Màu sắc không được là tín hiệu duy nhất; trạng thái phải có nhãn chữ hoặc biểu tượng đi kèm.

\subsection{Quy tắc đặt mockup trong báo cáo}

Mockup cấp phân hệ được đặt ngay trước các use case của phân hệ đó ở Phần~\ref{sec:phase2-by-uc}. Mỗi use case ghi rõ Screen ID mà actor sử dụng; một màn hình dùng chung chỉ xuất hiện một lần và các use case còn lại dẫn chiếu đến cùng hình. Cách tổ chức này giữ cho UI, sequence và activity của một nghiệp vụ nằm gần nhau nhưng không sao chép lại cùng một mockup.

\subsection{Working demonstration bằng chuỗi màn hình}

Phần này chọn một luồng xuyên suốt để trình bày prototype: \textit{Tra cứu Hub $\rightarrow$ xem tài nguyên $\rightarrow$ đặt xe $\rightarrow$ nhận xe $\rightarrow$ trả xe}. Mỗi bước phải trỏ được về Screen ID và Use Case ID tương ứng.

\ArtifactSection{CLICKABLE PROTOTYPE FLOW}{UI-DEMO-01}{Trình bày chuỗi màn hình hoặc liên kết prototype, kèm mô tả đầu vào, thao tác và kết quả tại từng bước.}{Luồng demonstration bằng chuỗi màn hình}
